\section{What has been established, and what has not}
\label{sec:status}

\subsection{Established}

Every statement in the following list is a theorem proved in the text or a result verified by the
programs of Appendix~\ref{app:data}; the label of each machine check is given in parentheses, and
Appendix~\ref{app:checks} lists all of them with their results. All
{{GEN:nchecks_rust}} checks of the program pass.
\begin{enumerate}
\item On the $7+1$ spacetime a single Dirac field has no Lorentz-scalar pair, because both
charge-conjugation matrices are symmetric; a doublet has exactly two, $\Phi^\pm_{U2}$, antisymmetric in
the doublet index; $\Phi^+_{U2}$ is the pair of total spin zero and orbital angular momentum zero; the
pairs join equal chiralities (R3-A1 to R3-A14).
\item The unitary transformations of the doublet index that preserve $\Phi^+_{U2}$ form $SU(2)$; with the
fermion number they form $U(2)_{U2}$. The covariant derivative~\eqref{eq:D} is covariant, its field
strength transforms by conjugation, and the pair is an $SU(2)_{U2}$ singlet of $U(1)_{U2}$ charge 1
(R3-A9, R3-G1 to R3-G6). The Dirac doublet is free of gauge anomalies; a chiral doublet cannot carry the
gauged $U(1)_{U2}$ (R3-B2).
\item The mean-field Dirac equation has the explicit spinor solutions~\eqref{eq:spinors} with the
quasi-particle energies~\eqref{eq:Epm}; the spin-zero pair opens the full gap at a Fermi surface, the
pseudoscalar pair only $\Delta_{U2}q^F_{U2}/\mu_{U2}$ (R3-C1 to R3-C6, R3-U1).
\item The effective potential~\eqref{eq:Omega} has a Mexican-hat minimum, which breaks
$U(2)_{U2}\to SU(2)_{U2}$, in the vacuum if and only if $\lambda^P_{U2}>\max(1,\lambda^S_{U2})$
(critical couplings $90\pi^4\Lambda^{-6}_{U2}$ and $288\pi^4\Lambda^{-6}_{U2}$ in the two
regularizations), and at finite density for every attraction, with the weak-coupling law~\eqref{eq:bcs};
a density can drive a first-order transition from the massive symmetric vacuum to the superfluid
(R3-P1 to R3-P12).
\item After the breaking: a Goldstone mode, absorbed by $B_{U2A}$, which acquires the mass
$g'_{U2}f_{U2}$; massless $SU(2)_{U2}$ bosons; and the neutral scalar $\sigma_{U2}$ with
$m_{\sigma U2}=2|\Delta_{U2}|$ exactly in the vacuum without scalar coupling, a bound state below the
pair-breaking threshold when the fermions have a bare mass, whose propagator is the pairing
interaction of the broken phase. At finite density the stiffness obeys Leggett's theorem and the
Goldstone velocity is $7^{-1/2}$ at weak coupling (R3-M1 to R3-M12).
\item No symbol of this paper denotes two objects, here or in the rest of the project, and the
documents are self-contained (R3-S1 to R3-S6 and the checks of self-containedness).
\end{enumerate}

\subsection{Not established}

\begin{enumerate}
\item \emph{Beyond mean field.} All results on the potential and the modes are of mean-field and
random-phase order. With two flavours there is no small parameter that controls the corrections. The
transition at the critical coupling should have the exponents of mean-field theory, because the
order parameter lives in eight dimensions, above the upper critical dimension four of a
Ginzburg--Landau theory; this has not been verified, and it would not make the mean-field values of
$\Delta_{U2}$ exact.
\item \emph{Renormalizability.} The theory is not renormalizable: $G^S_{U2}$ and $G^P_{U2}$ have the
dimension mass$^{-6}$, $g_{U2}$ and $g'_{U2}$ mass$^{-2}$, so every order of perturbation theory
brings higher powers of the cutoff. The cutoff $\Lambda_{U2}$ is part of the model, which is an
effective theory below it. The results that do not depend on the regularization are the existence of
a critical coupling, the Mexican hat, the Cooper instability, the Goldstone mode, $m_{\sigma U2}=2|\Delta_{U2}|$,
$m_{BU2}=g'_{U2}f_{U2}$, the massless $SU(2)_{U2}$ bosons, Leggett's theorem and $(v^G_{U2})^2=1/7$. The
numbers in units of $\Lambda_{U2}$ depend on it: the critical couplings differ by the factor $3.2$ between
(R7) and (R8), the vacuum gaps by the factor ${{GENM:num_gap_ratio_l105}}$ at $\lambda^P_{U2}=1.05$, and (R7) breaks Lorentz invariance
($f^s_{U2}\neq f^t_{U2}$). Whether an eight-dimensional theory of this kind can be made renormalizable
is not addressed here.
\item \emph{The gauge fields in the gap equation.} The exchange of $\boldsymbol W_{U2A}$ (attractive in
the singlet) and of $B_{U2A}$ (repulsive) is not included in the gap equation; the contact coupling
$G^P_{U2}$ stands for the net attraction. The couplings $g_{U2}$ and $g'_{U2}$ enter only through the
masses and the static potentials of Sec.~\ref{sec:breaking}.
\item \emph{Finite temperature.} Only zero temperature is computed; the critical temperature is not.
\item \emph{The neutralizing background} at finite density is assumed, not derived.
\item \emph{Global anomalies} of a chiral doublet are not examined; they do not concern the Dirac
doublet, which admits a gauge-invariant mass.
\item \emph{Physics in three space dimensions.} The model lives on the $7+1$ spacetime with momenta in
all seven space dimensions; how it would appear in a world with three large space dimensions is
not addressed, and gravity is not included.
\end{enumerate}

\subsection{Where the result differs from the wording of the task}

\begin{enumerate}
\item \emph{A ``scalar gauge boson'' does not exist.} A gauge boson is the quantum of a connection, a
one-form; its quanta are vectors. The task asks for an electrically neutral scalar boson of the new
group that appears only after the breaking and mediates the pairing. The model has two neutral
bosons that exist as particles only after the breaking: the massive vector $B_{U2A}$ (a gauge boson,
the analogue of the Z boson) and the scalar $\sigma_{U2}$ (not a gauge boson, the analogue of the Higgs
boson). The scalar is the one that carries the pairing interaction, in the sense of
Eq.~\eqref{eq:Vsigma}.
\item \emph{``A new group''.} As an abstract group, $U(2)_{U2}$ is the group of the electroweak theory.
What is new is its action on an internal index of 16-component fermions of the $7+1$ spacetime, the
reason for it (Theorem~1), and the field that breaks it (a pair of fermions).
\item \emph{``Mediated''.} The scalar mediates the pairing in the broken phase; the attraction that
forms the condensate in the first place is the contact coupling (or the exchange of the
$SU(2)_{U2}$ bosons).
\item \emph{Superfluid or superconductor.} The pair carries the gauged charge of $U(1)_{U2}$, so with
respect to $U(1)_{U2}$ the condensate is a superconductor (it expels the field of $B_{U2A}$); with respect
to electromagnetism, under which the doublet is assumed neutral, it is a superfluid.
\end{enumerate}
